\chapter{Números reales}\label{ch:b1:reals}

Todo el análisis descansa sobre una propiedad que distingue $\R$ de
$\Q$: todo \omterm{def:b1:logic:sets}{conjunto} no vacío y acotado superiormente tiene una
\emph{mínima} \omterm{def:b1:reals:bounds}{cota superior}. Este capítulo la enuncia con precisión,
deduce sus primeras consecuencias --- la \omterm{thm:b1:reals:archimedes}{propiedad arquimediana}, la
\omterm{thm:b1:reals:floor}{parte entera}, la densidad de los racionales y de los irracionales ---
y monta el vocabulario (sup, inf, máx, mín) que se usará
constantemente a partir del~\cref{ch:b1:seq}.

\section{La propiedad del supremo}

\begin{definition}[Cotas, supremo e ínfimo]\label{def:b1:reals:bounds}
Sea $A \subseteq \R$ no vacío. Un real $M$ es una \emph{cota
superior}\index{cota superior} de $A$ cuando $a \leq M$ para todo
$a \in A$; $A$ está \emph{acotado superiormente} cuando tiene una cota
superior (análogamente por abajo, con cotas inferiores;
\emph{acotado} significa las dos cosas). Un \emph{máximo} de $A$ es
una cota superior que pertenece a $A$.

El \emph{supremo} $\sup A$\index{supremo} es la mínima cota superior
de $A$, cuando existe; el \emph{ínfimo} $\inf A$\index{ínfimo} es la
máxima cota inferior.
\end{definition}

\begin{theorem}[Axioma de completitud de $\R$]\label{thm:b1:reals:sup}
$\R$ es un \omterm{def:b1:structures:field}{cuerpo} ordenado que contiene a $\Q$ y en el que \emph{todo
subconjunto no vacío y acotado superiormente tiene
\omterm{def:b1:reals:bounds}{supremo}}.\index{completitud de R@completitud de $\R$}
\end{theorem}
\admitted

\begin{remark}
Lo tomamos como axioma que define $\R$; construir un modelo (con
cortaduras de Dedekind o con sucesiones de Cauchy de racionales) y
demostrar su unicidad es honesto pero largo, y queda para estudios
posteriores. Obsérvese que $\Q$ no cumple la propiedad:
$\{x \in \Q : x^2 < 2\}$ está acotado superiormente pero no tiene
mínima \omterm{def:b1:reals:bounds}{cota superior} \emph{en $\Q$} --- su candidato, $\sqrt 2$, falta
(\cref{ex:b1:logic:sqrt2}). Pasando a los opuestos
($\sup(-A) = -\inf A$), todo \omterm{def:b1:logic:sets}{conjunto} no vacío y acotado inferiormente
tiene \omterm{def:b1:reals:bounds}{ínfimo}.
\end{remark}

\begin{proposition}[La caracterización con $\varepsilon$]\label{prop:b1:reals:epsilon}
Sean $A \neq \emptyset$ acotado superiormente y $s \in \R$. Entonces
$s = \sup A$ si y solo si
\begin{enumerate}
  \item $s$ es \omterm{def:b1:reals:bounds}{cota superior}: $\forall a \in A$, $a \leq s$; y
  \item nada menor lo es: $\forall \varepsilon > 0$, $\exists a \in
        A$, $a > s - \varepsilon$.
\end{enumerate}
\end{proposition}

\begin{proof}
Si $s = \sup A$: (1) se cumple por definición y, para (2),
$s - \varepsilon < s$ no es \omterm{def:b1:reals:bounds}{cota superior}, que es exactamente la
existencia de $a > s - \varepsilon$. Recíprocamente, (1) dice que $s$
es \omterm{def:b1:reals:bounds}{cota superior}; (2) dice que ningún $t < s$ es \omterm{def:b1:reals:bounds}{cota superior}
(tómese $\varepsilon = s - t$): $s$ es la mínima.
\end{proof}

\begin{example}\label{ex:b1:reals:supexamples}
$\sup \intoo{0}{1} = 1$, no alcanzado (no hay máximo);
$\sup \intcc{0}{1} = 1 = \max$. Para $A = \{1 - \frac 1n : n \in
\N^*\}$: $\sup A = 1$, no alcanzado; $\inf A = \min A = 0$. Un máximo,
cuando existe, es el \omterm{def:b1:reals:bounds}{supremo}; el sentido del $\sup$ es tener un
sustituto cuando el máximo no existe.
\end{example}

\begin{omfigure}
\begin{tikzpicture}[scale=8.5]
  \draw[-Stealth] (-0.05,0) -- (1.3,0);
  \fill[omDef] (0,0) circle (0.35pt) node[below=2pt, font=\small]
    {$0$};
  \fill[omDef] (0.5,0) circle (0.35pt) node[below=2pt, font=\small]
    {$\tfrac12$};
  \fill[omDef] (0.6667,0) circle (0.35pt)
    node[below=2pt, font=\small] {$\tfrac23$};
  \fill[omDef] (0.75,0) circle (0.35pt)
    node[below=2pt, font=\small] {$\tfrac34$};
  \fill[omDef] (0.8,0) circle (0.35pt);
  \fill[omDef] (0.8333,0) circle (0.35pt);
  \fill[omDef] (0.8571,0) circle (0.35pt);
  \fill[omDef] (0.875,0) circle (0.35pt);
  \fill[omDef] (0.8889,0) circle (0.35pt);
  \fill[omDef] (0.9,0) circle (0.35pt);
  \draw[omThm, thick] (1,-0.015) -- (1,0.015);
  \node[above, omThm, font=\small] at (1,0.02) {$\sup A = 1$};
  \draw[omProp, very thick] (1,-0.04) -- (1.27,-0.04);
  \node[below, omProp, font=\small] at (1.13,-0.05)
    {cotas superiores de $A$};
  \draw[Stealth-, gray] (0.93,0.03) -- (0.8,0.075)
    node[above left, font=\small, text=gray, inner sep=1pt]
    {algún $a > 1 - \varepsilon$};
\end{tikzpicture}

{\small El \omterm{def:b1:logic:sets}{conjunto} $A = \{1 - \frac1n : n \in \N^*\}$ sobre la recta
numérica: sus puntos se acumulan hacia $1$ sin alcanzarlo. Todo número
$\geq 1$ es \omterm{def:b1:reals:bounds}{cota superior} (la semirrecta) y nada menor lo es, porque
un elemento de $A$ entra en cada \omterm{prop:b1:reals:intervals}{intervalo}
$\intoo{1 - \varepsilon}{1}$: las dos cláusulas de la~\cref{prop:b1:reals:epsilon} en una sola
imagen. El \omterm{def:b1:reals:bounds}{supremo} es el extremo izquierdo de la semirrecta de
\omterm{def:b1:reals:bounds}{cotas superiores} --- y el axioma de completitud es precisamente la
garantía de que esa semirrecta tiene siempre extremo izquierdo.}
\end{omfigure}

\begin{example}[Calcular supremos en la práctica]\label{ex:b1:reals:supcompute}
Dos ejercicios completos con la~\cref{prop:b1:reals:epsilon}.

\emph{El \omterm{def:b1:logic:sets}{conjunto} $A = \{x + \frac1x : x > 0\}$.} Para todo $x > 0$,
$x + \frac1x - 2 = \frac{(\,\sqrt x - 1/\sqrt x\,)^2}{1} \geq 0$, luego
$2$ es cota inferior; y $2 = 1 + \frac11 \in A$: por tanto
$\inf A = \min A = 2$, alcanzado en $x = 1$. Por arriba, $A$ no está
acotado ($x + \frac1x > x$ puede superar cualquier $M$, por
el~\cref{thm:b1:reals:archimedes}): $\sup A$ no existe en $\R$ (vale
$+\infty$ en $\overline\R$).

\emph{El \omterm{def:b1:logic:sets}{conjunto} $B = \bigl\{\frac{m}{m + n} : m, n \in
\N^*\bigr\}$.} Todo elemento está en $\intoo{0}{1}$, luego $0$ y $1$
son cotas. Ninguna se alcanza: $\frac{m}{m+n} = 1$ obligaría a
$n = 0$. Para el \omterm{def:b1:reals:bounds}{supremo}, fíjese $n = 1$ y hágase crecer $m$:
$\frac{m}{m+1} = 1 - \frac{1}{m+1} > 1 - \varepsilon$ en cuanto
$m + 1 > \frac1\varepsilon$ (Arquímedes): $\sup B = 1$.
Simétricamente ($m = 1$, $n$ grande), $\inf B = 0$. La idea de cierre:
para fijar un \omterm{def:b1:reals:bounds}{supremo} basta un \emph{camino de un parámetro} bien
elegido dentro del \omterm{def:b1:logic:sets}{conjunto} --- aquí el camino $n = 1$ --- y la
caracterización con $\varepsilon$ no pide más.
\end{example}

\begin{example}[El espejo del ínfimo]\label{ex:b1:reals:infmirror}
El \omterm{def:b1:reals:bounds}{ínfimo} tiene su propia caracterización con $\varepsilon$,
obtenida de la~\cref{prop:b1:reals:epsilon} mediante
$\inf A = -\sup(-A)$: $i = \inf A$ si y solo si $i$ acota $A$
inferiormente y, para todo $\varepsilon > 0$, algún $a \in A$ cumple
$a < i + \varepsilon$. Un ejercicio con las dos cotas a la vez: sea
\[
A = \Bigl\{(-1)^n + \frac1n : n \in \N^*\Bigr\}
= \Bigl\{0,\ \tfrac32,\ -\tfrac23,\ \tfrac54,\ -\tfrac45,\
\dots\Bigr\} .
\]
Los índices pares dan $1 + \frac1n \leq \frac32$, con igualdad en
$n = 2$: como además los valores de índice impar son $\leq 0 <
\frac32$, resulta $\sup A = \max A = \frac32$. Los índices impares dan
$-1 + \frac1n > -1$, decreciendo hacia $-1$: todo elemento de $A$ es
$> -1$, y a $-1 + \varepsilon$ lo supera $-1 + \frac1n$ para impares
$n > \frac1\varepsilon$: $\inf A = -1$, no alcanzado. Un solo
\omterm{def:b1:logic:sets}{conjunto} y los cuatro comportamientos a la vista: un \omterm{def:b1:reals:bounds}{supremo} que es
máximo y un \omterm{def:b1:reals:bounds}{ínfimo} que no es mínimo.
\end{example}

\begin{remark}[Errores frecuentes con sup e inf]
Cuatro errores explican casi todos los puntos perdidos. (i)
\emph{Confundir $\sup$ y $\max$}: $\sup A$ no tiene por qué pertenecer
a $A$; escríbase $\max$ solo tras exhibir un elemento de $A$ que sea
\omterm{def:b1:reals:bounds}{cota superior}. (ii) \emph{Pasar desigualdades estrictas al
\omterm{def:b1:reals:bounds}{supremo}}: si $a < b$ para todo $a \in A$, solo se puede concluir
$\sup A \leq b$ --- testigo $A = \intoo{0}{1}$, $b = 1$. (iii)
\emph{Escribir $\sup A$ antes de comprobar que es lícito}: el símbolo
exige $A$ no vacío y acotado superiormente
(\cref{met:b1:reals:supproofs}); $\sup \emptyset$ y $\sup \N$ no están
definidos en $\R$ (los convenios de $\overline\R$ son un acto aparte y
explícito). (iv) \emph{Operaciones con \omterm{def:b1:logic:sets}{conjuntos}}:
$\sup(A \cup B) = \max(\sup A, \sup B)$ siempre, pero para $A \cap B$
no hay nada general --- puede ser vacío y, aun no siéndolo,
$\sup(A \cap B)$ puede quedar muy por debajo de $\min(\sup A, \sup B)$:
tómense $A = \{0, 2\}$ y $B = \{0, 3\}$, donde $\sup(A \cap B) = 0$.
\end{remark}

\begin{example}[Los conjuntos finitos tienen máximo --- un lema usado en silencio]\label{ex:b1:reals:finitemax}
Todo $F \subseteq \R$ finito y no vacío tiene máximo (y mínimo).
Inducción sobre el número de elementos: un \omterm{def:b1:logic:sets}{conjunto} unitario $\{a\}$
tiene $\max = a$; si la afirmación vale para \omterm{def:b1:logic:sets}{conjuntos} de $n$ elementos
y $F$ tiene $n + 1$, tómese cualquier $a \in F$: el \omterm{def:b1:logic:sets}{conjunto}
$F \setminus \{a\}$ tiene un máximo $m$, y $\max F$ es $m$ si
$a \leq m$, y $a$ en caso contrario. No interviene la completitud
---esto es puro orden más inducción, válido ya en $\Q$---, y sin
embargo el lema merece un \omterm{def:b1:logic:statement}{enunciado} honesto, porque las demostraciones
que vienen lo invocan en silencio: la construcción de la \omterm{thm:b1:reals:floor}{parte entera}
de más abajo («un \omterm{def:b1:logic:sets}{conjunto} de enteros atrapado en un rango finito tiene
un elemento máximo»), toda cota $\max(\abs{u_0}, \dots, \abs{u_{N-1}},
\dots)$ del~\cref{ch:b1:seq}, todo «tómese el mayor de los finitos
$\delta$» del~\cref{ch:b1:continuity}. Los \omterm{def:b1:logic:sets}{conjuntos} infinitos son donde
mueren los máximos y toman el relevo los \omterm{def:b1:reals:bounds}{supremos}: este capítulo existe
por el caso infinito.
\end{example}

\begin{theorem}[Propiedad arquimediana]\label{thm:b1:reals:archimedes}
Para todo $x \in \R$ existe $n \in \N$ con $n > x$. Equivalentemente:
para todos $\varepsilon > 0$ e $y > 0$, algún múltiplo $n\varepsilon$
supera a $y$.\index{propiedad arquimediana}
\end{theorem}

\begin{proof}
Supóngase que no: algún $x$ es \omterm{def:b1:reals:bounds}{cota superior} de $\N$. Entonces existe
$s = \sup \N$ (\cref{thm:b1:reals:sup}). Por
la~\cref{prop:b1:reals:epsilon}~(2) con $\varepsilon = 1$, hay
$n \in \N$ con $n > s - 1$; pero entonces $n + 1 \in \N$ y $n + 1 > s$,
en contra de que $s$ sea \omterm{def:b1:reals:bounds}{cota superior}. Para la segunda forma, sean
$\varepsilon > 0$ e $y > 0$: la primera forma aplicada a
$x = \frac{y}{\varepsilon}$ produce $n \in \N$ con
$n > \frac{y}{\varepsilon}$ y, multiplicando por $\varepsilon > 0$
(que conserva las desigualdades estrictas), $n\varepsilon > y$.
Recíprocamente, la segunda forma con $\varepsilon = 1$ e $y = x$
recupera la primera para $x > 0$, y $n = 1$ resuelve $x \leq 0$: los
dos \omterm{def:b1:logic:statement}{enunciados} son estrictamente equivalentes.
\end{proof}

\begin{example}[Arquímedes en acción]\label{ex:b1:reals:archimedeswork}
Tres usos inmediatos, necesarios constantemente más adelante. (i)
\emph{Ningún real positivo queda por debajo de todos los $\frac1n$}:
si $0 < \varepsilon$, tómese $n > \frac1\varepsilon$; entonces
$\frac1n < \varepsilon$. Dicho de otro modo, $\R$ no contiene
infinitesimales --- el informal «$\frac1n$ se hace arbitrariamente
pequeño» es exactamente este teorema. (ii) \emph{Umbrales explícitos}:
¿cuán grande debe ser $n$ para que $\frac{1}{n^2} \leq 10^{-6}$? Basta
$n \geq 10^3$ --- Arquímedes garantiza que tales $n$ existen, y el
álgebra los localiza. (iii) \emph{Las potencias superan cualquier
cota}: $2^n \geq n + 1$ (inducción), luego, para todo $M$, alguna
potencia de $2$ supera a $M$: el crecimiento geométrico usado para los
diádicos del~\cref{exo:b1:reals:8}. La idea de cierre: la
\omterm{thm:b1:reals:archimedes}{propiedad arquimediana} es la licencia que hay detrás de toda frase
del tipo «tómese $n$ suficientemente grande» --- de aquí en adelante
usaremos esa frase con libertad, y este ejemplo es su justificación de
una vez por todas.
\end{example}

\begin{theorem}[Parte entera]\label{thm:b1:reals:floor}
Para todo $x \in \R$ existe exactamente un entero, la \emph{parte
entera}\index{parte entera} $\lfloor x \rfloor$, con
\[
\lfloor x \rfloor \leq x < \lfloor x \rfloor + 1 .
\]
\end{theorem}

\begin{proof}
\emph{Existencia.} El \omterm{def:b1:logic:sets}{conjunto} $E = \{k \in \Z : k \leq x\}$ no es
vacío: por el~\cref{thm:b1:reals:archimedes} hay $m \in \N$ con
$m > -x$, y entonces $-m < x$, luego $-m \in E$. Está acotado
superiormente (por cualquier entero $n > x$, que existe por la misma
razón), de modo que, siendo un \omterm{def:b1:logic:sets}{conjunto} de enteros atrapado en el rango
finito $\intint{-m}{n}$, tiene un elemento máximo $k = \max E$.
Entonces $k \leq x$, y $k + 1 \notin E$ significa $x < k + 1$.

\emph{Unicidad.} Si $k$ y $k'$ cumplen los dos las desigualdades,
entonces $k \leq x < k' + 1$ da $k \leq k'$, y simétricamente
$k' \leq k$.
\end{proof}

\begin{example}[Partes enteras en la práctica]\label{ex:b1:reals:floorpractice}
$\lfloor 3.7 \rfloor = 3$, $\lfloor 5 \rfloor = 5$ y
$\lfloor -3.7 \rfloor = -4$: la \omterm{thm:b1:reals:floor}{parte entera} va hacia \emph{abajo},
no hacia $0$. Dos consecuencias de la unicidad
del~\cref{thm:b1:reals:floor} que usaremos en silencio. Primera, para
$n \in \Z$,
\[
\lfloor x + n \rfloor = \lfloor x \rfloor + n ,
\]
porque $\lfloor x \rfloor + n$ es un entero que cumple las dos
desigualdades que definen la de $x + n$ --- y solo un entero lo hace.
Segunda, $\lfloor \, \cdot \, \rfloor$ es no decreciente: si
$x \leq y$, entonces $\lfloor x \rfloor \leq x \leq y < \lfloor y
\rfloor + 1$, y un entero $< \lfloor y \rfloor + 1$ es
$\leq \lfloor y \rfloor$. Cuidado, en cambio: $\lfloor 2x \rfloor \neq
2\lfloor x \rfloor$ en general; $x = 0.6$ da $\lfloor 1.2 \rfloor = 1
\neq 0 = 2\lfloor 0.6 \rfloor$.

Lo que \emph{sí} es cierto es una identidad que conviene guardar (la de
Hermite, en su caso más simple): para todo real $x$,
\[
\lfloor x \rfloor + \Bigl\lfloor x + \frac12 \Bigr\rfloor
= \lfloor 2x \rfloor .
\]
Escríbase $x = \lfloor x\rfloor + u$ con $u \in \intco{0}{1}$ y
sepárense dos casos. Si $u < \frac12$: el miembro izquierdo es
$\lfloor x\rfloor + \lfloor x\rfloor = 2\lfloor x\rfloor$, y
$2x = 2\lfloor x\rfloor + 2u$ con $2u \in \intco{0}{1}$, luego el
derecho vale también $2\lfloor x\rfloor$. Si $u \geq \frac12$: el
izquierdo es $\lfloor x\rfloor + (\lfloor x\rfloor + 1)$, y
$2u \in \intco{1}{2}$ hace que el derecho valga
$2\lfloor x\rfloor + 1$. La idea de cierre:
$\lfloor x + \frac12\rfloor$ es el \emph{redondeo} de $x$ al entero
más próximo, de modo que la identidad dice que \omterm{thm:b1:reals:floor}{parte entera} más
redondeo es igual a la \omterm{thm:b1:reals:floor}{parte entera} del doble --- y la separación de
casos según la parte fraccionaria $u$ es la técnica estándar detrás de
toda identidad con \omterm{thm:b1:reals:floor}{partes enteras} (los
\cref{exo:b1:reals:2,exo:b1:reals:3} también funcionan así).
\end{example}

\begin{theorem}[Densidad de $\Q$ y de $\R \setminus \Q$]\label{thm:b1:reals:density}
Entre dos reales cualesquiera $x < y$ hay un racional y un
irracional.\index{densidad}
\end{theorem}

\begin{proof}
\emph{Un racional.} Por el~\cref{thm:b1:reals:archimedes}, tómese
$n \in \N^*$ con $n > \frac{1}{y - x}$, de modo que $ny - nx > 1$. Sea
$m = \lfloor nx \rfloor + 1$. Por un lado,
$nx < \lfloor nx \rfloor + 1 = m$ (\cref{thm:b1:reals:floor}); por
otro, $m = \lfloor nx \rfloor + 1 \leq nx + 1 < ny$. Dividiendo entre
$n$: $x < \frac mn < y$.

\emph{Un irracional.} Aplíquese el punto anterior al par
$x - \sqrt 2 < y - \sqrt 2$: algún racional $q$ queda entre ellos, y
entonces $q + \sqrt 2 \in \intoo{x}{y}$ es irracional (si
$q + \sqrt 2$ fuese racional, también lo sería $\sqrt 2$).
\end{proof}

\begin{example}[Ejecutar la demostración de la densidad]\label{ex:b1:reals:densityrun}
La demostración es un algoritmo; ejecutémoslo con $x = 1.414$ e
$y = \sqrt 2$. Como $1.4142^2 = 1.99996164 < 2$, se tiene
$\sqrt 2 > 1.4142$, luego $y - x > 0.0002$ y
$\frac{1}{y - x} < 5000$: la elección $n = 5000$ es lícita. Entonces
$nx = 7070$, luego $m = \lfloor 7070 \rfloor + 1 = 7071$, y el racional
producido es
\[
\frac{m}{n} = \frac{7071}{5000} = 1.4142,
\qquad
1.414 < 1.4142 < \sqrt 2 .
\]
La idea de cierre: la demostración solo necesita un $n$ ligeramente
mayor que $\frac{1}{y-x}$, y devuelve el primer múltiplo de $\frac 1n$
más allá de $x$. La densidad no es un milagro abstracto --- es la
división larga disfrazada, tema desarrollado con amplitud en el
problema del fin de semana (\cref{pb:b1:reals:1}).
\end{example}

\begin{remark}[Dónde se usa la completitud a continuación]
\Cref{thm:b1:reals:sup} es el único axioma no algebraico de este
libro, y todo teorema de existencia del análisis es ese axioma con
otro traje: el teorema de convergencia monótona
(\cref{ch:b1:seq}), el teorema de Bolzano--Weierstrass
(\cref{ch:b1:topology}), los teoremas del valor intermedio y de los
valores extremos (\cref{ch:b1:continuity}) y la definición misma de la
integral como \omterm{def:b1:reals:bounds}{supremo} de sumas inferiores
(\cref{ch:b1:integration}). El volumen del tercer año construye la
teoría de la medida y los espacios de Hilbert sobre ese mismo único
axioma. Cuando una demostración de los capítulos que vienen saque un
número real de la nada, búsquese el \omterm{def:b1:reals:bounds}{supremo} escondido.
\end{remark}

\begin{remark}[Entre lo discreto y lo denso]
$\Z$ y $\Q$ ocupan extremos opuestos dentro de $\R$: alrededor de cada
entero hay un hueco de longitud $1$ sin ningún otro (lo discreto ---
que es lo que hace que la \omterm{thm:b1:reals:floor}{parte entera} esté bien definida), mientras
que entre dos reales cualesquiera hay infinitos racionales (la
densidad). Sorprendentemente, para los \emph{\omterm{def:b1:structures:subgroup}{subgrupos} aditivos} de
$\R$ no hay nada intermedio: el~\cref{exo:b1:reals:9} demuestra que un
\omterm{def:b1:structures:subgroup}{subgrupo} así es, o bien de la forma $\alpha\Z$ (discreto), o bien
denso --- dicotomía que sostiene la densidad de $\{\sin n\}$ en
el~\cref{ch:b1:seq} y el monstruo constructivo
del~\cref{pb:b1:continuity:1}. Los \omterm{def:b1:logic:sets}{conjuntos} generales, claro, mezclan
libremente los dos comportamientos: $\Z \cup \Q\cap\intcc{0}{1}$ es
discreto a lo lejos y denso en el medio.
\end{remark}

\begin{method}[Demostrar igualdades con sup e inf]\label{met:b1:reals:supproofs}
Para demostrar $\sup A = s$: compruébese que $s$ acota $A$
superiormente y prodúzcase después, para cada $\varepsilon > 0$ (o
para una sucesión $\varepsilon = \frac 1n$), un elemento de $A$ mayor
que $s - \varepsilon$. Para comparar \omterm{def:b1:reals:bounds}{supremos}, úsense:
$A \subseteq B \implies \sup A \leq \sup B$; y, para todos $a, b$:
$\sup(A + B) = \sup A + \sup B$, donde $A + B = \{a + b\}$
(\cref{exo:b1:reals:5}). No se escriba nunca $\sup A$ antes de saber
que $A$ es no vacío y acotado superiormente.
\end{method}

\section{Intervalos}

\begin{proposition}[Caracterización de los intervalos]\label{prop:b1:reals:intervals}
Un subconjunto $I \subseteq \R$ es un \emph{intervalo}\index{intervalo}
(de los tipos familiares $\intoo{a}{b}$, $\intcc{a}{b}$,
$\intco{a}{b}$, $\intoc{a}{b}$, semirrectas, $\R$, $\emptyset$,
\omterm{def:b1:logic:sets}{conjuntos} unitarios) si y solo si es \emph{convexo}:
\[
\forall x, y \in I,\ \forall z \in \R, \quad
x \leq z \leq y \implies z \in I .
\]
\end{proposition}

\begin{proof}
Todos los tipos enumerados son claramente convexos. Recíprocamente,
sea $I$ convexo y no vacío. Póngase $a = \inf I$ si $I$ está acotado
inferiormente, y $a = -\infty$ en caso contrario; análogamente
$b = \sup I$ o $+\infty$. Afirmamos que $\intoo{a}{b} \subseteq I
\subseteq \intcc{a}{b}$ (con los convenios obvios en $\pm\infty$). La
segunda inclusión es la definición de cota. Para la primera, sea
$z \in \intoo{a}{b}$: como $z > a$, $z$ no es cota inferior (o bien
$a = -\infty$), luego algún $x \in I$ cumple $x < z$; análogamente,
algún $y \in I$ cumple $y > z$; la convexidad pone $z \in I$.

Queda leer el tipo a partir de la doble inclusión
$\intoo{a}{b} \subseteq I \subseteq \intcc{a}{b}$: los \omterm{def:b1:logic:sets}{conjuntos}
encajados entre un \omterm{prop:b1:reals:intervals}{intervalo} abierto y su cierre difieren de
$\intoo{a}{b}$ solo por la presencia o ausencia de los extremos
(finitos). Explícitamente: si $a, b \in \R$, las cuatro posibilidades
para $(a \in I,\ b \in I)$ dan $\intoo{a}{b}$, $\intco{a}{b}$,
$\intoc{a}{b}$, $\intcc{a}{b}$ (incluidos los casos degenerados
$a = b$: \omterm{def:b1:logic:sets}{conjunto} unitario si $a \in I$); si $a = -\infty$ y
$b \in \R$, se obtiene $\intoo{-\infty}{b}$ o $\intoc{-\infty}{b}$;
simétricamente para $a \in \R$, $b = +\infty$; y $a = -\infty$,
$b = +\infty$ da $I = \R$. Todos los casos están en la lista: hecho.
\end{proof}

\begin{remark}[Por qué la convexidad es el test adecuado]
La proposición convierte una definición \emph{geométrica} (una lista
de diez formas) en un test \emph{lógico de una línea}, y el test es lo
que realmente se usa: para demostrar que un \omterm{def:b1:logic:sets}{conjunto} es un \omterm{prop:b1:reals:intervals}{intervalo},
no se persiga cuál de las diez formas es --- verifíquese la convexidad
y déjese que la proposición decida el tipo. El teorema del valor
intermedio del~\cref{ch:b1:continuity} se enunciará exactamente así
(«la imagen continua de un \omterm{prop:b1:reals:intervals}{intervalo} es un \omterm{prop:b1:reals:intervals}{intervalo}»), y su
demostración produce la convexidad, no la forma.
\end{remark}

\begin{remark}[Recta real ampliada]
Resulta cómodo añadir dos símbolos y trabajar en
$\overline\R = \R \cup \{-\infty, +\infty\}$, con los convenios
$\sup A = +\infty$ cuando $A$ no está acotado superiormente y
$\sup \emptyset = -\infty$. Entonces \emph{todo} subconjunto de $\R$
tiene \omterm{def:b1:reals:bounds}{supremo} en $\overline\R$ --- una comodidad de notación usada
con libertad para los límites del~\cref{ch:b1:seq}.
\end{remark}

\begin{example}[Calcular en $\overline\R$]\label{ex:b1:reals:extendedcompute}
Con los convenios en vigor: $\sup \Z = +\infty$, $\inf \Z = -\infty$;
para $A = \{n + (-1)^n n : n \in \N\} = \{0, 4, 0, 8, \dots\} \cup
\{0\}$, $\sup A = +\infty$ (los términos pares $2n$ no están acotados)
e $\inf A = \min A = 0$; y $\sup\emptyset = -\infty \leq \inf\emptyset
= +\infty$ --- el único \omterm{def:b1:logic:sets}{conjunto} cuyo \omterm{def:b1:reals:bounds}{supremo} es \emph{menor} que su
\omterm{def:b1:reals:bounds}{ínfimo}, recordatorio de que los convenios se eligen para que $\sup$
sea creciente e $\inf$ decreciente respecto de la inclusión:
\[
A \subseteq B \implies \sup A \leq \sup B
\quad\text{and}\quad \inf A \geq \inf B ,
\]
ahora válidas sin la salvedad de la no vacuidad. Lo que los convenios
\emph{no} proporcionan es aritmética: $+\infty + (-\infty)$ y
$0 \times (+\infty)$ quedan sin definir, y toda manipulación
algebraica de \omterm{def:b1:reals:bounds}{supremos} debe comprobar antes que no forma ninguna de
ellas. La recta ampliada es contabilidad, no un sistema numérico.
\end{example}

\begin{example}[El supremo que se escapó de $\Q$]\label{ex:b1:reals:supescape}
Volvamos al \omterm{def:b1:logic:sets}{conjunto} de la observación inicial,
$A = \{x \in \Q : x^2 < 2\}$, y calculemos su \omterm{def:b1:reals:bounds}{supremo} \emph{en $\R$}.
Es no vacío ($1 \in A$) y está acotado superiormente por $1.5$ (si
$x > 1.5$, entonces $x^2 > 2.25 > 2$), luego existe $s = \sup A$.
Afirmamos que $s = \sqrt 2$ (el número real construido en
el~\cref{exo:b1:reals:12}). \omterm{def:b1:reals:bounds}{Cota superior}: todo $a \in A$ cumple
$a < \sqrt2$ --- para $a \leq 0$ es evidente y, para $a > 0$,
$a \geq \sqrt2$ daría $a^2 \geq 2$. Nada menor sirve: dado
$t < \sqrt2$, la densidad (\cref{thm:b1:reals:density}) proporciona un
racional $q$ con $\max(1, t) < q < \sqrt 2$, y entonces $q^2 < 2$,
luego $q \in A$ supera a $t$. Por la~\cref{prop:b1:reals:epsilon},
$s = \sqrt2 \notin \Q$. La idea de cierre: el \omterm{def:b1:reals:bounds}{supremo} de un \omterm{def:b1:logic:sets}{conjunto}
de racionales no tiene por qué ser racional --- la completitud es
precisamente la promesa de que $\R$, a diferencia de $\Q$, no deja
escapar nunca un \omterm{def:b1:reals:bounds}{supremo}; este ejemplo es la observación inicial del
capítulo, ahora demostrada en lugar de solo señalada.
\end{example}

\begin{remark}[Perspectivas dentro de este volumen]
Las tres herramientas del capítulo tienen carreras distintas por
delante. El \omterm{def:b1:reals:bounds}{supremo} dirige la mitad analítica: los límites monótonos
(\cref{ch:b1:seq}), la definición misma de la integral
(\cref{ch:b1:integration}) y, en la geometría del~\cref{ch:b1:euclid},
la distancia de un punto a un subespacio --- un \omterm{def:b1:reals:bounds}{ínfimo} que la
proyección ortogonal convierte en mínimo. La \omterm{thm:b1:reals:floor}{parte entera} vuelve allí
donde lo discreto se encuentra con lo continuo: desarrollos en cifras
(el problema del fin de semana de este capítulo), la aproximación de
Dirichlet por el palomar (\cref{pb:b1:derivative:1}), las
comparaciones integrales de sumas (\cref{ch:b1:series}). Los
argumentos de densidad ascienden a método en
el~\cref{ch:b1:continuity}: una identidad entre funciones continuas
solo hay que comprobarla en $\Q$ --- la mitad de la ecuación funcional
de Cauchy (\cref{pb:b1:continuity:1}) es exactamente ese movimiento.
En caso de duda sobre de dónde saca una demostración de este volumen
sus \omterm{def:b1:logic:statement}{enunciados} de existencia, la respuesta es casi siempre: de este
capítulo.
\end{remark}

\section{Ejercicios}

\begin{exercise}[$\star$]\label{exo:b1:reals:1}
Determínense (con demostración) sup, inf, máx y mín --- cuando existan --- de:
\[
A = \Bigl\{\frac{1}{n} : n \in \N^*\Bigr\},
\qquad
B = \Bigl\{\frac{(-1)^n n}{n+1} : n \in \N\Bigr\},
\qquad
C = \{x \in \R : x^2 < 3\}.
\]
\end{exercise}

\begin{exercise}[$\star$]\label{exo:b1:reals:2}
Demuéstrese que, para todos $x, y \in \R$, $\lfloor x \rfloor +
\lfloor y \rfloor \leq \lfloor x + y \rfloor \leq \lfloor x \rfloor +
\lfloor y \rfloor + 1$, y que las dos cotas se alcanzan.
\end{exercise}

\begin{exercise}[$\star$]\label{exo:b1:reals:3}
Demuéstrese que, para todo $x \in \R$ y todo $n \in \N^*$,
$\Bigl\lfloor \frac{\lfloor nx \rfloor}{n} \Bigr\rfloor = \lfloor x
\rfloor$.
\end{exercise}

\begin{exercise}[$\star$]\label{exo:b1:reals:4}
Sean $A \subseteq B$ subconjuntos no vacíos de $\R$, con $B$ acotado.
Demuéstrese que $\inf B \leq \inf A \leq \sup A \leq \sup B$.
\end{exercise}

\begin{exercise}[$\star\star$]\label{exo:b1:reals:5}
Para $A, B \subseteq \R$ no vacíos y acotados, defínanse $A + B = \{a + b : a
\in A,\ b \in B\}$ y $-A = \{-a : a \in A\}$. Demuéstrese:
\[
\sup(A + B) = \sup A + \sup B,
\qquad
\sup(-A) = -\inf A .
\]
\end{exercise}

\begin{exercise}[$\star\star$]\label{exo:b1:reals:6}
Sean $f, g \colon E \to \R$ funciones acotadas. Demuéstrese que
\[
\sup_{x \in E}\, \bigl(f(x) + g(x)\bigr) \leq \sup_{x \in E} f(x) +
\sup_{x \in E} g(x),
\]
y dese un ejemplo en el que la desigualdad sea estricta. ¿Por qué esto
no contradice el~\cref{exo:b1:reals:5}?
\end{exercise}

\begin{exercise}[$\star\star$]\label{exo:b1:reals:7}
Demuéstrese que $\sqrt 2 + \sqrt 3$ es irracional. \emph{(Elévese al
cuadrado y úsese la irracionalidad de $\sqrt 6$, que se demuestra con
el~\cref{exo:b1:arith:7}.)}
\end{exercise}

\begin{exercise}[$\star\star$]\label{exo:b1:reals:8}
Demuéstrese que el \omterm{def:b1:logic:sets}{conjunto} $D = \bigl\{\frac{m}{2^n} : m \in \Z,\ n \in
\N\bigr\}$ de los racionales diádicos es denso en $\R$: entre dos
reales cualesquiera hay un racional diádico.
\end{exercise}

\begin{exercise}[$\star\star\star$]\label{exo:b1:reals:9}
Sea $G$ un \omterm{def:b1:structures:subgroup}{subgrupo} de $(\R, +)$ con $G \neq \{0\}$. Póngase
$\alpha = \inf\,(G \cap \intoo{0}{+\infty})$. Demuéstrese:
\begin{enumerate}
  \item si $\alpha > 0$, entonces $G = \alpha\Z$;
  \item si $\alpha = 0$, entonces $G$ es denso en $\R$.
\end{enumerate}
Dedúzcase que $\Z + \sqrt 2\,\Z$ es denso en $\R$.
\end{exercise}

\begin{exercise}[$\star\star\star$]\label{exo:b1:reals:10}
Para $A, B$ \omterm{def:b1:logic:sets}{conjuntos} no vacíos de reales positivos, sea
$AB = \{ab : a \in A, b \in B\}$. Demuéstrese que
$\sup(AB) = \sup A \cdot \sup B$ (caso acotado), y véase con un ejemplo
que la positividad es esencial.
\end{exercise}

\begin{exercise}[$\star\star$]\label{exo:b1:reals:11}
Para $A \subseteq \R$ no vacío y acotado, defínase el
\emph{diámetro}
\[
\operatorname{diam} A = \sup\,\{\abs{a - a'} : a, a' \in A\} .
\]
Demuéstrese que $\operatorname{diam} A = \sup A - \inf A$ y que
$\intcc{\inf A}{\sup A}$ es el menor \omterm{prop:b1:reals:intervals}{intervalo} cerrado que contiene
a $A$.
\end{exercise}

\begin{exercise}[$\star\star\star$]\label{exo:b1:reals:12}
Sean $y > 0$ y $E = \{x \geq 0 : x^2 \leq y\}$. Demuéstrese que $E$ es
no vacío y acotado superiormente, y que $s = \sup E$ cumple $s^2 = y$
\emph{(descártense $s^2 < y$ y $s^2 > y$ exhibiendo, en cada caso, un
$h > 0$ pequeño que contradiga la definición del \omterm{def:b1:reals:bounds}{supremo})}. Dedúzcase
que todo $y > 0$ tiene una única raíz cuadrada $\sqrt y > 0$ y que
$y \mapsto \sqrt y$ es creciente en $\intoo{0}{+\infty}$.
\end{exercise}

\section{Problema: desarrollos en cifras y el ritmo de los racionales}

\begin{problem}[{Problema del fin de semana --- desarrollos $b$-ádicos:
existencia, unicidad y la periodicidad que caracteriza $\Q$}]
\label{pb:b1:reals:1}
Todo real de $\intco{0}{1}$ tiene un desarrollo en cifras en toda base
$b \geq 2$; el desarrollo es único una vez prohibidas las colas
formadas por la cifra $b - 1$; y es finalmente periódico exactamente
cuando el número es racional. Este problema demuestra los tres hechos
solo con el axioma de completitud --- sin sucesiones, sin series: solo
el \omterm{def:b1:reals:bounds}{supremo}, la \omterm{thm:b1:reals:archimedes}{propiedad arquimediana} y la \omterm{thm:b1:reals:floor}{parte entera} --- y cierra
con el argumento diagonal de Cantor en forma de cifras. En todo el
problema, $b \geq 2$ es un entero fijo (la \emph{base}), una
\emph{cifra} es un elemento de $\intint{0}{b-1}$, y una cadena de
cifras $(d_n)_{n \geq 1}$ es \emph{propia} cuando no es finalmente
igual a $b - 1$ (es decir: para todo $N$ existe $n > N$ con
$d_n \leq b - 2$).\index{desarrollo b-adico@desarrollo
$b$-ádico}\index{desarrollo decimal}

\medskip
\noindent\textbf{Parte I --- Cifras a mano.} División larga de $p$
entre $q$ en base $b$: multiplíquese el resto actual por $b$, divídase
entre $q$, anótese el cociente como cifra siguiente y guárdese el resto.
\begin{enumerate}
  \item En base $10$, ejecútese el algoritmo con $\frac 18$ y con
        $\frac 17$, anotando en cada paso la cifra \emph{y} el resto.
        Compruébese que los restos de $\frac 17$ recorren el ciclo
        $1, 3, 2, 6, 4, 5$ y que las cifras $142857$ se repiten después
        para siempre.
  \item Calcúlense los desarrollos en base $2$ de $\frac 13$ y de
        $\frac{5}{16}$, y el desarrollo en base $3$ de $\frac 12$.
        Obsérvese: un número termina y los otros dos se repiten --- y
        $\frac 12$, tan dócil en base $10$, se repite para siempre en
        base $3$.
  \item Para $x = \frac pq \in \intco{0}{1}$ irreducible, véase que las
        cifras que produce el algoritmo son finalmente todas $0$ si y
        solo si el resto $b^N p \bmod q$ se anula para algún $N$, si y
        solo si $q$ \omterm{def:b1:arith:divides}{divide} a alguna potencia $b^N$, si y solo si todo
        factor \omterm{def:b1:arith:prime}{primo} de $q$ \omterm{def:b1:arith:divides}{divide} a $b$. Compruébese:
        $\frac{1}{20}$ termina en base $10$, pero no en base $3$.
  \item Defínase el \emph{truncamiento}
        $s_n = \lfloor b^n x \rfloor / b^n$. Para $x = \sqrt 2$ y
        $b = 10$, calcúlense $s_0, \dots, s_4$ verificando en cada paso
        que dos cuadrados consecutivos encierran a $2$ (por ejemplo,
        $1.4142^2 = 1.99996164 < 2 < 2.00024449 = 1.4143^2$), y
        compruébese cada vez $s_n \leq \sqrt 2 < s_n + 10^{-n}$.
\end{enumerate}

\medskip
\noindent\textbf{Parte II --- Existencia, a partir del \omterm{def:b1:reals:bounds}{supremo}.}
Fíjese $x \in \intco{0}{1}$ y póngase $A_n = \lfloor b^n x \rfloor$ y
$d_n = A_n - b\,A_{n-1}$ para $n \geq 1$.
\begin{enumerate}[resume]
  \item Véase que $A_0 = 0$ y $b\,A_{n-1} \leq A_n \leq b\,A_{n-1} +
        b - 1$; conclúyase que cada $d_n$ es una cifra.
  \item Véase que $s_n := A_n b^{-n}$ cumple
        \[
        s_n = \sum_{k=1}^{n} d_k\,b^{-k}
        \qquad\text{and}\qquad
        s_n \leq x < s_n + b^{-n} .
        \]
  \item Demuéstrese $b^n \geq n + 1$ por inducción y véase después que
        $(s_n)$ es no decreciente y que $x = \sup_n s_n$
        \emph{(úsense la~\cref{prop:b1:reals:epsilon} y
        el~\cref{thm:b1:reals:archimedes})}.
  \item Véase que la cadena $(d_n)$ es propia: si $d_k = b - 1$ para
        todo $k > N$, calcúlese $s_n$ para $n > N$ con una suma
        geométrica finita y contradígase la pregunta 6.
  \item Recíprocamente, sea $(e_n)_{n \geq 1}$ una cadena propia de
        cifras cualquiera y $t_n = \sum_{k=1}^n e_k b^{-k}$. Véase que
        $y = \sup_n t_n$ existe, está en $\intco{0}{1}$ y cumple
        $t_n \leq y < t_n + b^{-n}$ para todo $n$ \emph{(para la
        desigualdad estricta, úsese una cifra $e_m \leq b - 2$ con
        $m > n$)}. Dedúzcase que $\lfloor b^n y \rfloor = b^n t_n$ y,
        después, que las cifras de $y$, en el sentido de la
        pregunta 5, son exactamente los $e_n$.
\end{enumerate}

\medskip
\noindent\textbf{Parte III --- Unicidad, orden, desplazamiento.}
\begin{enumerate}[resume]
  \item Reúnanse las preguntas 5--9 en el \emph{teorema del desarrollo
        $b$-ádico}: las aplicaciones $x \mapsto (d_n)$ y
        $(e_n) \mapsto \sup_n t_n$ son biyecciones mutuamente inversas
        entre $\intco{0}{1}$ y el \omterm{def:b1:logic:sets}{conjunto} de las cadenas propias de
        cifras. En particular, dos cadenas propias distintas no tienen
        nunca el mismo valor.
  \item Permítanse ahora las cadenas impropias. Véase que una cadena
        con $e_n = b - 1$ para todo $n > M$ (con $M \geq 0$ mínimo)
        tiene valor $t_M + b^{-M}$; conclúyase que
        $0.999\dots = 1$ en base $10$ y que los reales con dos
        representaciones en cifras son exactamente las fracciones
        $b$-ádicas $m/b^N \in \intoo{0}{1}$ --- cualquier otro real
        tiene una sola, incluso entre las cadenas impropias.
  \item Demuéstrese que la \omterm{def:b1:logic:inj}{biyección} de la pregunta 10 conserva el
        orden lexicográfico: si las cadenas propias de $x$ y de $y$
        difieren por primera vez en el índice $m$, entonces $x < y$ si
        y solo si $d_m < e_m$.
  \item (Lema del desplazamiento) Sea $x \in \intco{0}{1}$ de cifras
        $(d_n)$. Véase que la parte fraccionaria de $bx$ tiene cifras
        $(d_{n+1})_{n \geq 1}$ \emph{(calcúlese
        $\lfloor b^n(bx - A_1)\rfloor$ usando $\lfloor u - K \rfloor =
        \lfloor u \rfloor - K$ para $K$ entero)}, y dedúzcase por
        inducción que la parte fraccionaria de $b^m x$ tiene cifras
        $(d_{n+m})_{n \geq 1}$.
\end{enumerate}

\medskip
\noindent\textbf{Parte IV --- La racionalidad es periodicidad.} Sea
$x = \frac pq \in \intco{0}{1}$ irreducible y $r_n = b^n p \bmod q$ el
resto de la división euclídea de $b^n p$ entre $q$.
\begin{enumerate}[resume]
  \item Véase que $A_n = \dfrac{b^n p - r_n}{q}$ y que
        $r_n = (b\,r_{n-1}) \bmod q$.
  \item Véase que $d_n = \Bigl\lfloor \dfrac{b\,r_{n-1}}{q}
        \Bigr\rfloor$: cada cifra es función solo del resto anterior.
        Esta es exactamente la división larga de la parte I.
  \item Aplíquese el principio del palomar
        (\cref{cor:b1:counting:pigeonhole}) a $r_0, \dots, r_q$ y
        conclúyase: el desarrollo de todo racional es finalmente
        periódico, con anteperíodo y período a lo sumo $q$.
  \item Recíprocamente, supóngase que las cifras de $y \in \intco{0}{1}$
        son \emph{puramente} periódicas: $d_{n+T} = d_n$ para todo
        $n \geq 1$. Usando el lema del desplazamiento y la unicidad de
        la pregunta 10, véase que la parte fraccionaria de $b^T y$ es
        igual a $y$ y dedúzcase que $(b^T - 1)\,y \in \N$: así pues,
        $y$ es racional con denominador que \omterm{def:b1:arith:divides}{divide} a $b^T - 1$.
        Verifíquese el mecanismo en $0.(142857)$:
        $142857 \times 7 = 999999$.
  \item Trátese el caso finalmente periódico desplazando, y enúnciese
        el \emph{criterio de periodicidad}: $x \in \intco{0}{1}$ es
        racional si y solo si su desarrollo $b$-ádico propio es
        finalmente periódico --- en una base si y solo si en todas.
  \item Para $x = \frac 1q$ con $\gcd(q, b) = 1$, véase que el
        desarrollo es puramente periódico y que su período mínimo es el
        menor $T \geq 1$ con $b^T \equiv 1 \pmod q$ (el orden
        multiplicativo de $b$ \omterm{def:b1:complex:field}{módulo} $q$). Compruébese que, para
        $q = 7$, $b = 10$, las potencias de $10$ \omterm{def:b1:complex:field}{módulo} $7$ recorren
        $3, 2, 6, 4, 5, 1$: orden $6$, de acuerdo con la pregunta 1.
\end{enumerate}

\medskip
\noindent\textbf{Parte V --- Dividendos y la diagonal.}
\begin{enumerate}[resume]
  \item Sea $x^*$ el real de $\intco{0}{1}$ cuyas cifras en base $10$
        valen $1$ en las posiciones triangulares $\frac{j(j + 1)}{2}$
        ($j \geq 1$) y $0$ en las demás: $x^* =
        0.101001000100001\dots$ Véase que su cadena de cifras es propia
        pero no finalmente periódica \emph{(un período $T$ obligaría a
        unos separados por huecos de a lo sumo $T$, pero los huecos
        crecen)}, y conclúyase que $x^*$ es irracional: un número
        demostrado irracional por puro ritmo.
  \item Véase que, para toda base $b$, el \omterm{def:b1:logic:sets}{conjunto}
        $\{m/b^n : m \in \Z, n \in \N\}$ es denso en $\R$
        (generalizando el~\cref{exo:b1:reals:8}), y que todo racional
        $\frac pq \in \intoo{0}{1}$ tiene un desarrollo
        \emph{finito} en base $q$. Moraleja: ser finito es una
        propiedad del par (número, base); la periodicidad --- la
        racionalidad --- es intrínseca.
  \item (Diagonal de Cantor) Sea $k \mapsto x_k$ una \omterm{def:b1:logic:map}{aplicación}
        cualquiera de $\N^*$ en $\intco{0}{1}$. Defínase la cadena de
        cifras $e_k = 1$ si la $k$-ésima cifra de $x_k$ es distinta de
        $1$, y $e_k = 2$ en caso contrario. Véase que $(e_k)$ es
        propia, que su valor $y$ está en $\intco{0}{1}$ y que
        $y \neq x_k$ para todo $k$. Conclúyase: ninguna \omterm{def:b1:logic:map}{aplicación}
        $\N^* \to \intco{0}{1}$ es \omterm{def:b1:logic:inj}{sobreyectiva}. (El vocabulario de la
        numerabilidad, y el hogar propio de este teorema, es
        el~\cref{ch:b1:topology}.)
  \item Véase que si los desarrollos propios de $x$ y de $y$ coinciden
        hasta el índice $n$, entonces $\abs{x - y} < b^{-n}$, y
        refútese el recíproco con $x = 0.1$, $y = 0.0999$ en base $10$:
        la proximidad de los números no obliga a que coincidan las
        cifras. ¿Qué reales tienen la culpa?
  \item Ejecútese la parte IV con $x = \frac{1}{10}$ en base $b = 2$:
        calcúlense restos y cifras hasta que se repitan, y conclúyase
        $\frac{1}{10} = (0.0\overline{0011})_2$, con anteperíodo $1$ y
        período $4$. Explíquese, con la pregunta 3, por qué ninguna
        cadena binaria finita será jamás igual a $\frac{1}{10}$ --- la
        razón por la que el $0.1 + 0.2$ en coma flotante de un
        ordenador no vale exactamente $0.3$.
  \item Síntesis. En una frase cada uno: ¿dónde ha usado la
        demostración (i) la completitud, (ii) la
        \omterm{thm:b1:reals:archimedes}{propiedad arquimediana}, (iii) la cláusula de unicidad de la
        \omterm{thm:b1:reals:floor}{parte entera}, (iv) el principio del palomar? Y la moraleja:
        $\intco{0}{1}$ queda fielmente codificado por las cadenas
        propias de cifras y la racionalidad se lee como periodicidad
        --- y, sin embargo, el análisis prefiere el \omterm{def:b1:reals:bounds}{supremo} a las
        cifras. ¿Por qué? (Piénsese en sumar dos cadenas de cifras.)
\end{enumerate}
\end{problem}
